\documentclass[10pt,a4paper]{amsart}
\usepackage[margin=19mm]{geometry}
\usepackage{amsmath,amssymb,mathtools}
\usepackage[scr=rsfs,cal=euler]{mathalfa}
\usepackage{xfrac}
\usepackage[unicode,hidelinks,bookmarks=false]{hyperref}
\newcommand{\ie}{i.e.\ }
\newcommand{\cf}{cf.\ }
\newcommand{\resp}{respectively\ }
\newcommand{\Subsection}{\subsection}
\providecommand{\nobreakdash}{\nobreak\hbox{-}}
\setlength{\emergencystretch}{2em}
\multlinegap0pt
\usepackage{amsaddr}
\newtheorem{theorem}{Theorem}
\newtheorem{lemma}{Lemma}
\newtheorem{corollary}{Corollary}
\newtheorem{remark}{Remark}
\newtheorem{definition}{Definition}
\newcommand{\norm}[1]{\left\lVert#1\right\rVert}
\DeclareMathOperator{\Ima}{Im}
\newcommand{\LD}[1]{{\color{blue} #1}}
\newcommand{\ND}[1]{{\color{red} #1}}
\begin{document}
\title[Dirichlet-Neumann waveform relaxation for heterogeneous heat]{Dirichlet-Neumann waveform relaxation for heterogeneous heat equations: continuous and time discrete L2 analysis}
\author{Niklas Kotarsky, Philipp Birken, Martin J. Gander, Lu-di Lu}
\date{}
\maketitle
\noindent\emph{Source and licence.} Dirichlet-Neumann waveform relaxation for heterogeneous heat equations: continuous and time discrete L2 analysis. Version arXiv:2606.29987v1 (CC BY 4.0). This material is distributed under the Creative Commons Attribution 4.0 International licence, http://creativecommons.org/licenses/by/4.0/, as recorded in the arXiv OAI licence metadata for this exact version and shown on the abstract page https://arxiv.org/abs/2606.29987v1. Full rights notice: licenses/arxiv-2606.29987-CC-BY-4.0.txt.\par\smallskip
\noindent\textit{Editorial scope.} Counted unit: the original \texttt{theorem} environment \texttt{thm:timeDisc-case2} at source lines 554--571 together with its complete proof at lines 574--582; the delivered document keeps source lines 92--602 so that every referenced label and every citation resolves inside it. Native editable source is the paper's own concatenated TeX; the AMS wrapper is editorial formatting only. No publisher class, upstream script or unrelated artwork is redistributed.
\par\smallskip\noindent\textit{Reference note.} This article ships no compiled bibliography (biblatex); the entries delivered here are composed from the author's own BibTeX (.bib) fields, with no field invented and no entry dropped.
\section{Continuous convergence analysis}\label{sec:2}
	
We consider as our model problem two coupled linear heat equations in one dimension. Let $Q_1 := (-a, 0) \times (0, T)$ and $Q_2 := (0, b) \times (0, T)$ denote two space-time domains with $a>0$, $b>0$, and the time $T>0$. In each domain $Q_j$, we denote by $u_j$ the temperature, $f_j$ the external force, $u_{0, j}$ the initial condition, $\alpha_j$ the volumetric heat capacity, and $\lambda_i$ the heat conductivity. For each domain, the associated heat equation reads
\begin{equation}\label{eq:CoupledSystem}
	\begin{aligned}
	\alpha_1 \partial_t u_1 -\lambda_1 \partial_{xx} u_1 &= f_1& &\text{in } Q_1,&
	\alpha_2 \partial_t u_2 -\lambda_2 \partial_{xx} u_2 &= f_2&  &\text{in } Q_2,\\
	u_1(-a, t) &= g_l(t)& &\text{in } (0,T),&	
	u_1(b, t) &= g_r(t)&  &\text{in } (0,T),\\
	u_1(x, 0) &= u_{0, 1}(x)& &\text{in } (-a,0),& 
	u_2(x, 0) &= u_{0, 2}(x)& &\text{in } (0,b),
\end{aligned}
\end{equation}
where $g_l$ and $g_r$ are two Dirichlet boundary conditions. These two heat equations are coupled at the interface using the two coupling conditions
\begin{equation}\label{eq:CoupleCondition}
u_1(0, t) = u_2(0, t) \ \text{ in } (0,T), \quad 
-\lambda_1 \partial_x u_1(0, t) = -\lambda_2 \partial_x u_2(0, t) \ \text{ in } (0,T).
\end{equation}
The first condition imposes continuity of temperature, whereas the second ensures continuity of the heat flux at the interface.


The goal of our study is to analyze convergence of the Dirichlet--Neumann waveform relaxation (DNWR) method applied to our model problem~\eqref{eq:CoupledSystem}-\eqref{eq:CoupleCondition}. For any given initial guess $h^0(t)$ and the iteration index $k = 1, 2, \ldots$, the DNWR iteration reads
\begin{equation}\label{eq:DNWR}
\begin{aligned}
	\alpha_1 \partial_t u_1^{k+1} - \lambda_1 \partial_{xx} u_1^{k+1} &= f_1 \ \text{ in } Q_1,&
	\alpha_2 \partial_t u_2^{k+1} - \lambda_2 \partial_{xx} u_2^{k+1} &= f_2 \ \text{ in } Q_2, \\
	u_1^{k+1}(-a, t) &= g_l(t), \ &
	-\lambda_1 \partial_x u_1^{k+1}(0, t) &= -\lambda_2 \partial_x u_2^{k+1}(0, t), \\
	u_1^{k+1}(0, t) &= h^{k}(t), &
	u_1^k(b, t) &= g_r(t), \\
	u_1^{k+1}(x, 0) &= u_{0, 1}(x), &
	u_2^{k+1}(x, 0) &= u_{0, 2}(x),
\end{aligned}
\tag{DNWR}
\end{equation}
and we update the Dirichlet transmission condition $h^k$ by
\[h^{k+1}(t) := (1-\theta) h^k(t) + \theta u^{k+1}_2(0, t),
\quad
\theta \in (0, 1),\]
where $\theta$ is the relaxation parameter. Note that here we choose to impose the Dirichlet condition at the interface in $Q_1$ and the Neumann condition at the interface in $Q_2$. 

To present our convergence analysis, we first introduce a framework based on the exponentially weighted Fourier transform proposed in~\cite{Kwok2025}.


\begin{definition}\label{def:FourierTransform}
Let $r>0$, $\omega\in\mathbb{R}$, and $\phi$ a given function. The exponentially weighted Fourier transform of the function $\phi$ is defined by
\[\hat{\phi}(r + i\omega) :=  \int_{-\infty}^{\infty} \phi(t) e^{-(r+i \omega) t} \, dt,\] 
if this integral is finite.
\end{definition}
Note that the exponentially weighted Fourier transform of a function $\phi(t)$ is the standard Fourier transform applied to the function $e^{-rt}\phi(t)$, with $r>0$ the weight parameter. We use the hat notation to denote such a transform, and write the argument as $r+i\omega$ to emphasize the dependence on both the weight parameter $r$ and the frequency $\omega$. From Definition~\ref{def:FourierTransform}, we can verify that the exponentially weighted Fourier transform inherits most of the properties of the Fourier transform, such as linearity, time shifting, and differentiation. Also, it exists if $\phi\in L^1$. 
The Fourier transform of an $L^2$ function is in general defined via extension using the Plancherel theorem, see e.g.~\cite{grafakos2014a}. For the exponentially weighted Fourier transform of an $L^2$ function, we need additional assumptions, since the weight $e^{-rt}$ blows up as $t\to-\infty$ for $r>0$. The next result establishes the existence of this transform under an assumption that is relevant to our study.


\begin{lemma}\label{prepCont2norm}
Let $r > 0$ and $\phi$ be a function defined in $(-\infty,\infty)$ such that $\norm{\phi}_{L^2(\mathbb{R^+})} < \infty$ and $\phi(t) = 0$ for $t < 0$. Then the exponentially weighted Fourier transform of $\phi$ exists, and we have
\[\frac{1}{2 \pi} \int_{-\infty}^{\infty} |\hat{\phi}(r+i\omega)|^2 d\omega \leq \int_{-\infty}^{\infty} |\phi(t)|^2 dt.\]
\end{lemma}

\begin{proof}
For $r > 0$ and $\phi(t) = 0$ for $t < 0$, we have
\[\int_{-\infty}^{\infty} \big(e^{-rt} \phi(t)\big)^2 \, dt =  \int_{0}^{\infty} \big(e^{-rt}\phi(t)\big)^2 \, dt \leq \int_{0}^{\infty} \big(\phi(t)\big)^2 \, dt < \infty.\]
Thus, $\psi(t) := e^{-r t} \phi(t)$ is in $L^2(\mathbb{R})$. Using the result in~\cite{grafakos2014a}, the Fourier transform of $\psi$ is well defined, meaning that the exponentially weighted Fourier transform of $\phi$ is also well defined. 
Applying now the Plancherel theorem to $\psi$, we obtain 
\[\frac{1}{2 \pi} \int_{-\infty}^{\infty} |\hat{\phi}(r+i\omega)|^2 \, d\omega = \int_{-\infty}^{\infty} \big(e^{-rt} \phi(t)\big)^2 \, dt \leq \int_{0}^{\infty} \big(\phi(t)\big)^2 \, dt = \int_{-\infty}^{\infty} \big(\phi(t)\big)^2 \, dt.\]
\end{proof}


To analyze convergence, we apply the exponentially weighted Fourier transform to the error equation associated with~\eqref{eq:DNWR} and derive the contraction factor. The error equation is obtained by subtracting the solution from the model problem~\eqref{eq:CoupledSystem}-\eqref{eq:CoupleCondition}, or equivalently, by setting directly $f_1$, $f_2$, $u_{0,1}$, $u_{0,2}$, $g_l$, and $g_r$ all to zero in~\eqref{eq:DNWR}. Before applying the exponentially weighted Fourier transform to~\eqref{eq:DNWR}, we also need to extend the initial guess $h^0$ to the time interval $(-\infty, \infty)$ by setting it to zero outside its original domain. The iterates $(u_1^k, u_2^k, h^k)$ for $k>0$ are then extended to $(-\infty,\infty)$ through the iteration itself. For brevity, we retain the notation $h^0$ and $(u_1^k, u_2^k, h^k)$ for the extended functions. 


For $\omega\in\mathbb{R}$ and $r\in\mathbb{R}^+$, the exponentially weighted Fourier transform applied to the error equation associated with~\eqref{eq:DNWR} gives
\begin{equation}\label{eq:transformedDNWR}
\begin{aligned}
	\alpha_1 (r+i \omega) \hat{u}_1^{k+1} -\lambda_1 \partial_{xx} \hat{u}_1^{k+1} &= 0 \quad \text{ in } (-a,0),\\
	\hat{u}_1^{k+1}(-a, r+i \omega) &= 0,\\
	\hat{u}_1^{k+1}(0, r+i\omega) &= \hat h^{k}(r+i \omega),\\
	\alpha_2 (r+i \omega) \hat{u}_2^{k+1} -\lambda_2 \partial_{xx} \hat{u}_2^{k+1} &= 0 \quad \text{ in } (0,b),\\
	-\lambda_1 \partial_x \hat{u}_1^{k+1}(0, r+i \omega) &= -\lambda_2 \partial_x \hat{u}_2^{k+1}(0, r+i \omega),\\
	u_1^k(b, r+i \omega) &= 0, \\
	\hat{h}^{k+1}(r+i\omega) = (1-\theta) \hat{h}^k&(r+i\omega) + \theta \hat{u}^{k+1}_2(0, r+i \omega).
\end{aligned}
\end{equation}
These are two coupled second-order ordinary differential equations in space, and solving them for given $\omega$ and $r$ yields the closed form solutions
\[\begin{aligned}
	\hat{u}_1^{k+1}(x, r+i \omega) &= \frac{\sinh\left( \sqrt{\frac{\alpha_1 (r+i \omega)}{\lambda_1}} (x +a) \right)}{\sinh\left( \sqrt{\frac{\alpha_1 (r+i \omega)}{\lambda_1}} a \right)} \hat{h}^k(r+i \omega), \\
	\hat{u}_2^{k+1}(x, r+i \omega) &= 
	\sqrt{\frac{\alpha_1 \lambda_1}{\alpha_2 \lambda_2}}  \coth\left( \sqrt{\frac{\alpha_1 (r+i \omega)}{\lambda_1}} a\right) \\
	&\times \frac{\sinh\left( \sqrt{\frac{\alpha_2 (r+i \omega)}{\lambda_2}} (x - b) \right)}{\cosh\left( \sqrt{\frac{\alpha_2 (r+i \omega)}{\lambda_2}} b \right)} \hat{h}^k(r+i \omega).
\end{aligned}\]
Using the last equation in~\eqref{eq:transformedDNWR}, we obtain $\hat{h}^k(r+i \omega) = \hat{\rho}(r+i \omega)^k \hat{h}^0(r+i \omega)$, where $\hat{\rho}$ is given by
\begin{equation}\label{eq:cvfCont}
\hat{\rho}(s) := 1 - \theta - \theta \sqrt{\frac{\alpha_1 \lambda_1}{\alpha_2 \lambda_2}}\coth\left(\sqrt{\frac{\alpha_1 s}{\lambda_1}} a\right) \tanh\left(\sqrt{\frac{\alpha_2 s}{\lambda_2}} b\right), \quad s\in\mathbb{C}.
\end{equation}
The function $\hat{\rho}$ evaluated at $s=r+i\omega$ with $\text{Re}(s) = r>0$ is the contraction factor associated with the transformed DNWR iteration~\eqref{eq:transformedDNWR}. Note that $ \hat{\rho}(r+i \omega)$ depends on the domain size $a$ and $b$, and the physical parameters $\alpha_j$ and $\lambda_j$, $j=1, 2$. 
Regarding the relaxation parameter $\theta$, one wants the associated contraction factor to be small, such that ~\eqref{eq:DNWR} converges very fast. Before discussing the choice of $\theta$, we first present an error estimate for arbitrary relaxation parameter.

\begin{theorem}\label{2normSuperCont}
	For any given $\theta\in (0, 1)$ and $h^0\in L^2(0, T)$, the error of the iteration~\eqref{eq:DNWR} in the time interval $(0, T)$ satisfies the estimate 
	\begin{equation}\label{eq:GenEstCon}
	\|h^k\|_{L^2(0, T)} \leq \min_{r > 0} e^{rT}\max_{\omega \in \mathbb{R}} \left|\hat{\rho}(r + i \omega)\right|^k \, \norm{h^0}_{L^2(0, T)}.
	\end{equation}
\end{theorem} 

\begin{proof}
	By \cite[Chapter 4]{Lions72}, $h^k$ exists and is well defined for $h^0 \in L^2(0,T)$. For $r > 0$, we have
	\[\begin{aligned}
		\|h^k\|_{L^2(0, T)}^2 = \int_{0}^{T} \left(h^{k}(t)\right)^2 dt = \int_{0}^{T} e^{2rt} e^{-2 r t} \left(h^{k}(t)\right)^2 &dt \leq e^{2rT} \int_{0}^{T} \left(e^{-r t}h^{k}(t)\right)^2 dt \\
		&\leq e^{2rT} \int_{-\infty}^{\infty} \left( e^{- r t} h^{k}(t)\right)^2 dt,
	\end{aligned}\]
	where in the last inequality we use the  extension of $h^k$ to the interval $(-\infty,\infty)$. By Lemma~\ref{prepCont2norm}, the exponentially weighted Fourier transform of $h^k$ exists. Applying the Plancherel theorem, we find
	\[\begin{aligned}
		e^{2rT} \int_{-\infty}^{\infty} \left( e^{- r t}h^{k}(t)\right)^2 dt 
		&= \frac{e^{2rT}}{2 \pi} \int_{-\infty}^{\infty} \left| \hat{h}^{k}(r + i \omega) \right|^2 d\omega \\
		&= \frac{e^{2rT}}{2 \pi} \int_{-\infty}^{\infty} \left| \hat{\rho}(r + i \omega)^k \hat{h}^0(r + i \omega) \right|^2 d\omega \\
		&\leq \frac{e^{2rT}}{2 \pi} \max_{\omega \in \mathbb{R}} \left|\hat{\rho}(r  + i\omega)\right|^{2k} \int_{-\infty}^{\infty} \left| \hat{h}^{0}(r + i \omega) \right|^2 d\omega \\
		&= e^{2rT} \max_{\omega \in \mathbb{R}} \left|\hat{\rho}(r + i\omega)\right|^{2k}\int_{-\infty}^{\infty} \left| h^{0}(t) \right|^2 dt \\
		&= e^{2rT} \max_{\omega \in \mathbb{R}} \left|\hat{\rho}(r + i\omega)\right|^{2k}\int_{0}^{T} \left| h^{0}(t) \right|^2 dt,
	\end{aligned}\]
	where we use Lemma~\ref{prepCont2norm} once again, and the fact that $h^0$ is extended from $(0, T)$ to $(-\infty, \infty)$ by 0 for the last equality. Combining these computations gives 
	\[\|h^k\|_{L^2(0, T)}^2 \leq e^{2rT} \max_{\omega \in \mathbb{R}} |\hat{\rho}(r + i\omega)|^{2k} \|h^0\|_{L^2(0, T)}^2.\]
	Taking the square root on both sides and minimizing over all $r>0$ leads to~\eqref{eq:GenEstCon}.
\end{proof}


Theorem~\ref{2normSuperCont} provides a general error estimate of~\eqref{eq:DNWR} for any relaxation parameter. To achieve fast convergence, we can choose a relaxation parameter $\theta$ such that the convergence factor~\eqref{eq:cvfCont} is small. In particular, the choice
\begin{equation}\label{eq:thtopt}
\theta = \frac1{1 + \sqrt{\frac{\alpha_1 \lambda_1}{\alpha_2 \lambda_2}}\coth\left(\sqrt{\frac{\alpha_1 (r+i \omega)}{\lambda_1}} a\right) \tanh\left(\sqrt{\frac{\alpha_2 (r+i \omega)}{\lambda_2}} b\right)}
\end{equation}
yields convergence in two iterations. However, this choice depends both on the weight parameter $r$ and the Fourier frequency $\omega$, making it not useful in practice. On the other hand, when the physics in both subdomains satisfy $\sqrt{\alpha_1/\lambda_1} a = \sqrt{\alpha_2/\lambda_2} b$, then the contraction factor~\eqref{eq:cvfCont} reduces to the simple form $\hat{\rho} = 1 - \theta - \theta\sqrt{\alpha_1 \lambda_1}/\sqrt{\alpha_2 \lambda_2}$, which is independent of the frequency $\omega$ and the weight parameter $r$. In this case, we obtain the optimal relaxation parameter
\begin{equation}\label{eq:thetaopt}
	\theta^* := \frac{\sqrt{\alpha_2\lambda_2}}{\sqrt{\alpha_1\lambda_1}+\sqrt{\alpha_2\lambda_2}}.
\end{equation}
Note that the choice $\theta$ given in~\eqref{eq:thtopt} converges to $\theta^*$ when $\omega$ goes to infinity.
Therefore, for high frequencies, the DNWR iteration always converges very fast for the choice of relaxation parameter $\theta^*$, and thus the convergence will be independent of the mesh size, as discussed in~\cite[Chapter 4.7, p.174]{Ciaramella2022}. 
Furthermore, the parameter $\theta^*$ is much more convenient to use in practice, as it only depends on the physics of the problem. Therefore, we are interested in the impact of $\theta^*$ on the convergence of~\eqref{eq:DNWR} in more general cases, {\it i.e.}, $\sqrt{\alpha_1/\lambda_1} a \neq \sqrt{\alpha_2/\lambda_2} b$. Substituting $\theta^*$ into ~\eqref{eq:cvfCont} and using properties of the differences of arguments for hyperbolic functions, we obtain
\begin{equation}\label{eq:rhotheta}
\hat{\rho}(s)|_{\theta^*} = -\frac{1}{1+\sqrt{\frac{\alpha_2\lambda_2}{\alpha_1\lambda_1}}}\frac{\sinh\left(\left(\sqrt{\frac{\alpha_1}{\lambda_1}} a - \sqrt{\frac{\alpha_2}{\lambda_2}} b\right) \sqrt{s}\right)}{\sinh\left(\sqrt{\frac{\alpha_1}{\lambda_1}} a\sqrt{s}\right)\cosh\left(\sqrt{\frac{\alpha_2}{\lambda_2}} b\sqrt{s}\right)},
\quad 
s \in \mathbb{C}.
\end{equation}
To simplify the notation, we also introduce the function \begin{equation}\label{def:H}
	\hat H(s) := \frac{\sinh\left(|\alpha-\beta|\sqrt{s}\right)}{\sinh\left(\alpha\sqrt{s}\right)\cosh\left(\beta\sqrt{s}\right)}, \quad s \in \mathbb{C}.
\end{equation}
With the help of~\eqref{def:H}, the error estimate~\eqref{eq:GenEstCon} associated with $\theta = \theta^*$ is then
\begin{equation}\label{eq:SpeEstCon}
\|h^k\|_{L^2(0, T)} \leq \left(\frac{1}{1+\sqrt{\frac{\alpha_2\lambda_2}{\alpha_1\lambda_1}}}\right)^k  \min_{r > 0} e^{rT}\max_{\omega \in \mathbb{R}} \left|\hat H(r + i \omega)\right|^k \, \norm{h^0}_{L^2(0, T)},
\end{equation}
where $\alpha = \sqrt{\alpha_1/\lambda_1} a$ and $\beta = \sqrt{\alpha_2/\lambda_2} b$ in $\hat H$ for \eqref{def:H}. Therefore, it is sufficient to study the min-max problem associated with $\hat H$. We first solve the maximization problem with respect to the frequency $\omega$, then use $r$ as an auxiliary parameter to find a superlinear estimate. The next result shows the global maximum of the maximization problem.


\begin{lemma}\label{lem:max}
	Let $\alpha, \beta, r>0$ be given. Then the function $\hat H$ defined in~\eqref{def:H} satisfies
	$\max_{\omega \in \mathbb{R}}|\hat H(r + i\omega)| = \hat H(r)$.
\end{lemma}
\begin{proof}
	In the proof of~\cite[Theorem 2]{ma:13}, it is shown that the function $\hat H(s)$ with $\alpha >0$, $\beta > 0$ and $\text{Re}(s)>0$, is the Laplace transform of a positive function denoted by $H(t)$. Thus, we have
\[\begin{aligned}
\max_{\omega \in \mathbb{R}} |\hat{H}(r + i \omega)| =  \max_{\omega \in \mathbb{R}} \Big| \int_{-\infty}^{\infty} &H(t) e^{-r t - i \omega t} dt \Big|
\leq \max_{\omega \in \mathbb{R}} \int_{-\infty}^{\infty} |H(t)|e^{-rt} dt  \\
&= \int_{-\infty}^{\infty} H(t)e^{-rt} dt  = \hat{H}(r).
\end{aligned}\] 
To see that this is actually an equality, it suffices to substitute $\omega = 0$ into $\hat{H}(r + i \omega)$, which gives $\hat{H}(r)$, and thus the proof is complete.
\end{proof}


Applying Lemma~\ref{lem:max} to \eqref{eq:SpeEstCon}, we have
\begin{equation}\label{eq:min-cont}
\norm{h^k}_{L^2(0, T)} \leq \left(\frac{1}{1+\sqrt{\frac{\alpha_2\lambda_2}{\alpha_1\lambda_1}}}\right)^k\min_{r>0}e^{rT} \left(\hat H(r)\right)^k \norm{h^0}_{L^2(0, T)},
\end{equation}
It remains to treat the minimization problem related to $r$. Depending on the physics of the problem, we have two results.


\begin{theorem}[Case $\sqrt{\alpha_1/\lambda_1} a > \sqrt{\alpha_2/\lambda_2} b$]\label{thm:cont-case1}
If $\theta = \theta^*$ and $\sqrt{\alpha_1/\lambda_1} a > \sqrt{\alpha_2/\lambda_2} b$, then the error of the algorithm~\eqref{eq:DNWR} satisfies 
\begin{equation}\label{eq:ContEst1}
\norm{h^k}_{L^2(0, T)} \leq \left(2\frac{1 - \sqrt{\frac{\alpha_2 \lambda_1}{\alpha_1\lambda_2}}\frac{b}{a}}{1 + \sqrt{\frac{\alpha_2 \lambda_2}{\alpha_1\lambda_1}}} \right)^k e^{-\frac{\alpha_2 k^2b^2}{4\lambda_2 T}} \norm{h^0}_{L^2(0, T)}.
\end{equation}
\end{theorem}


\begin{proof}
The idea is to first bound the function $\hat H(r)$ by an exponential function, and then to analyze the exponential function to find a minimum.

Let $0<\alpha<\beta$ and $f(x) = \beta\sinh(\alpha x) - \alpha \sinh(\beta x)$. We have $f'(x) = \alpha\beta(\cosh(\alpha x) - \cosh(\beta x)) < 0$, i.e., $f$ is decreasing. Hence $f(x) \leq f(0) = 0$, $\forall x>0$, yielding $\sinh(\alpha x) / \sinh(\beta x) \leq \alpha / \beta$. Taking $\alpha = \sqrt{\alpha_1/\lambda_1} a - \sqrt{\alpha_2/\lambda_2} b$, $\beta = \sqrt{\alpha_1/\lambda_1} a$ and $x = \sqrt{r}$, we find
\begin{equation}\label{eq:sinhbound}
\frac{\sinh\left((\sqrt{\frac{\alpha_1}{\lambda_1}} a - \sqrt{\frac{\alpha_2}{\lambda_2}} b) \sqrt{r} \right)}{\sinh\left(\sqrt{\frac{\alpha_1}{\lambda_1}} a \sqrt{r}\right)}\leq \frac{\sqrt{\frac{\alpha_1}{\lambda_1}} a - \sqrt{\frac{\alpha_2}{\lambda_2}} b}{\sqrt{\frac{\alpha_1}{\lambda_1}} a} = 1 - \sqrt{\frac{\alpha_2 \lambda_1}{\alpha_1\lambda_2}}\frac{b}{a}.
\end{equation}
On the other hand, using the definition of the hyperbolic cosine, we have
\begin{equation}\label{eq:coshbound}
\frac{1}{\cosh\left(\sqrt{\frac{\alpha_2}{\lambda_2}} b \sqrt{r} \right)} 
= \frac{2 \exp\left(-\sqrt{\frac{\alpha_2}{\lambda_2}} b \sqrt{r}\right)}{1 + \exp \left(-2\sqrt{\frac{\alpha_2}{\lambda_2}} b \sqrt{r} \right)} 
\leq 2 \exp\left(-\sqrt{\frac{\alpha_2}{\lambda_2}} b\sqrt{r}\right).
\end{equation}
The product of~\eqref{eq:sinhbound} and~\eqref{eq:coshbound} to the power $k$ gives
\begin{equation}\label{eq:BoundThm2}
\left(\frac{\sinh\left(\left(\sqrt{\frac{\alpha_1}{\lambda_1}} a - \sqrt{\frac{\alpha_2}{\lambda_2}} b\right) \sqrt{r} \right)}{\sinh\left(\sqrt{\frac{\alpha_1}{\lambda_1}} a \sqrt{r}\right)\cosh\left(\sqrt{\frac{\alpha_2}{\lambda_2}} b \sqrt{r}\right)}\right)^k
\leq 2^k\left(1 - \sqrt{\frac{\alpha_2 \lambda_1}{\alpha_1\lambda_2}}\frac{b}{a}\right)^k \exp\left(-kb\sqrt{\frac{\alpha_2r}{\lambda_2}}\right).
\end{equation}
Therefore, the estimate~\eqref{eq:min-cont} with the definition~\eqref{def:H} of $\hat H$ becomes
\begin{equation}\label{eq:minBoundThm2}
\begin{aligned}
\norm{h^k}_{L^2(0, T)} \leq& \left(\frac{1}{1+\sqrt{\frac{\alpha_2\lambda_2}{\alpha_1\lambda_1}}}\right)^k\min_{r>0}e^{rT} \left(\hat H(r)\right)^k \norm{h^0}_{L^2(0, T)}
\\
\leq& \left(2\frac{1 - \sqrt{\frac{\alpha_2 \lambda_1}{\alpha_1\lambda_2}}\frac{b}{a}}{1 + \sqrt{\frac{\alpha_2 \lambda_2}{\alpha_1\lambda_1}}} \right)^k \min_{r>0}e^{rT - kb\sqrt{\frac{\alpha_2r}{\lambda_2}}}\norm{h^0}_{L^2(0, T)}.
\end{aligned}
\end{equation}

It remains now to find the minimum of $\exp(rT - kb\sqrt{\alpha_2r}/\sqrt{\lambda_2})$. Differentiating it with respect to $r$ gives $(T - k b\sqrt{\alpha_2}/\sqrt{4\lambda_2 r}) \exp(rT - kb \sqrt{\alpha_2 r}/\sqrt{\lambda_2})$. Setting the derivative to 0 yields the extremal point
\begin{equation}\label{eq:ropt-case1}
r^*_{1c} = \frac{\alpha_2k^2 b^2}{4\lambda_2T^2}.
\end{equation}
Note that the derivative changes sign from negative to positive at $r^*_{1c}$, and thus $r^*_{1c}$ is a minimum. Substituting $r^*_{1c}$ into the second inequality in~\eqref{eq:minBoundThm2} gives the estimate~\eqref{eq:ContEst1}.
\end{proof}


\begin{theorem}[Case $\sqrt{\alpha_1/\lambda_1} a < \sqrt{\alpha_2/\lambda_2} b$]\label{thm:cont-case2}
Let $\theta = \theta^*$. If $\sqrt{\alpha_1/\lambda_1} a < \sqrt{\alpha_2/\lambda_2} b < 2\sqrt{\alpha_1/\lambda_1} a$, then the error of the algorithm~\eqref{eq:DNWR} satisfies 
\begin{equation}\label{eq:estimate-case2-1}
\norm{h^k}_{L^2(0, T)} \leq \left(2\frac{\sqrt{\frac{\alpha_2 \lambda_1}{\alpha_1\lambda_2}}\frac{b}{a} - 1}{\sqrt{\frac{\alpha_2 \lambda_2}{\alpha_1\lambda_1}}+1} \right)^k e^{-\frac{\alpha_2 k^2b^2}{4\lambda_2 T}} \norm{h^0}_{L^2(0, T)}.
\end{equation}
If $\sqrt{\alpha_2/\lambda_2} b \geq 2\sqrt{\alpha_1/\lambda_1} a$, then the error of the algorithm~\eqref{eq:DNWR} satisfies 
\begin{equation}\label{eq:estimate-case2-2}
\norm{h^k}_{L^2(0, T)} \leq \left(2\frac{\sqrt{\frac{\alpha_2 \lambda_1}{\alpha_1\lambda_2}}\frac{b}{a} - 1}{\sqrt{\frac{\alpha_2 \lambda_2}{\alpha_1\lambda_1}}+1} \right)^k e^{-\frac{\alpha_1 k^2a^2}{\lambda_1 T}} \norm{h^0}_{L^2(0, T)}.
\end{equation}
\end{theorem}

\begin{proof}
The idea is the same as in the proof of Theorem~\ref{thm:cont-case1}.
%\[\frac{\sinh\left(\left(\sqrt{\frac{\alpha_2}{\lambda_2}} b - \sqrt{\frac{\alpha_1}{\lambda_1}} a \right) \sqrt{r} \right)}{\sinh\left(\sqrt{\frac{\alpha_1}{\lambda_1}} a \sqrt{r}\right)\cosh\left(\sqrt{\frac{\alpha_2}{\lambda_2}} b \sqrt{r}\right)}.\]
If $\sqrt{\alpha_1/\lambda_1} a < \sqrt{\alpha_2/\lambda_2} b < 2\sqrt{\alpha_1/\lambda_1} a$, then $0<\sqrt{\alpha_2/\lambda_2} b - \sqrt{\alpha_1/\lambda_1} a < \sqrt{\alpha_1/\lambda_1} a$. We can thus follow the same steps as in the proof of Theorem~\ref{thm:cont-case1} by taking $\alpha = \sqrt{\alpha_2/\lambda_2} b - \sqrt{\alpha_1/\lambda_1} a$ and $\beta = \sqrt{\alpha_1/\lambda_1} a$.
%\[\sinh\left(\left(\sqrt{\frac{\alpha_2}{\lambda_2}} b - \sqrt{\frac{\alpha_1}{\lambda_1}} a \right) \sqrt{r} \right) \leq \sinh\left(\sqrt{\frac{\alpha_1}{\lambda_1}} a \sqrt{r}\right), \]
We find a similar bound as~\eqref{eq:BoundThm2}, namely
\begin{equation}\label{eq:BoundThm3}
	\left(\frac{\sinh\left(\left(\sqrt{\frac{\alpha_2}{\lambda_2}} b - \sqrt{\frac{\alpha_1}{\lambda_1}} a \right) \sqrt{r} \right)}{\sinh\left(\sqrt{\frac{\alpha_1}{\lambda_1}} a \sqrt{r}\right)\cosh\left(\sqrt{\frac{\alpha_2}{\lambda_2}} b \sqrt{r}\right)}\right)^k
	\leq 2^k\left(\sqrt{\frac{\alpha_2 \lambda_1}{\alpha_1\lambda_2}}\frac{b}{a} - 1\right)^k \exp\left(-kb\sqrt{\frac{\alpha_2r}{\lambda_2}}\right).
\end{equation}
This leads to the same miminizer $r^*_{1c}$ given in~\eqref{eq:ropt-case1}, and we obtain the estimate~\eqref{eq:estimate-case2-1}.

If $\sqrt{\alpha_2/\lambda_2} b \geq 2\sqrt{\alpha_1/\lambda_1} a$, we substitute $\tau=\exp(-2\sqrt{r})$ and obtain 
\[\frac{\sinh\left(\left(\sqrt{\frac{\alpha_2}{\lambda_2}} b - \sqrt{\frac{\alpha_1}{\lambda_1}} a \right) \sqrt{r} \right)}{\sinh\left(\sqrt{\frac{\alpha_1}{\lambda_1}} a \sqrt{r}\right)\cosh\left(\sqrt{\frac{\alpha_2}{\lambda_2}} b \sqrt{r}\right)}
= 2 \tau^{\sqrt{\frac{\alpha_1}{\lambda_1}} a}\frac{1-\tau^{\sqrt{\frac{\alpha_2}{\lambda_2}} b-\sqrt{\frac{\alpha_1}{\lambda_1}} a}}{\left(1+\tau^{\sqrt{\frac{\alpha_2}{\lambda_2}} b}\right)\left(1-\tau^{\sqrt{\frac{\alpha_1}{\lambda_1}} a}\right)}.\] 
The latter fraction can be bounded by
\[\frac{1-\tau^{\sqrt{\frac{\alpha_2}{\lambda_2}} b-\sqrt{\frac{\alpha_1}{\lambda_1}} a}}{\left(1+\tau^{\sqrt{\frac{\alpha_2}{\lambda_2}} b}\right)\left(1-\tau^{\sqrt{\frac{\alpha_1}{\lambda_1}} a}\right)} 
\leq \frac{1-\tau^{\sqrt{\frac{\alpha_2}{\lambda_2}} b-\sqrt{\frac{\alpha_1}{\lambda_1}} a}}{1-\tau^{\sqrt{\frac{\alpha_1}{\lambda_1}} a}}.\] 
Let now $\alpha \geq \beta >0$ and $f(x) = \beta(1 - e^{-\alpha x}) - \alpha(1 - e^{-\beta x})$, then $f'(x) = \alpha\beta(e^{-\alpha x} - e^{-\beta x}) \leq 0$, i.e., $f$ is non-increasing. Hence, $f(x) \leq f(0) = 0$, $\forall x > 0$, yielding $(1 - e^{-\alpha x})/((1 - e^{-\beta x})) \leq \alpha/\beta$. Taking then $\alpha = \sqrt{\alpha_2/\lambda_2} b - \sqrt{\alpha_1/\lambda_1} a$, $\beta = \sqrt{\alpha_1/\lambda_1} a$ and $x = 2\sqrt{r}$ gives
\begin{equation}\label{eq:bound-exp}
	\frac{1-\tau^{\sqrt{\frac{\alpha_2}{\lambda_2}} b-\sqrt{\frac{\alpha_1}{\lambda_1}} a}}{1-\tau^{\sqrt{\frac{\alpha_1}{\lambda_1}} a}} 
	\leq \frac{\sqrt{\frac{\alpha_2}{\lambda_2}} b-\sqrt{\frac{\alpha_1}{\lambda_1}} a}{\sqrt{\frac{\alpha_1}{\lambda_1}} a} = \sqrt{\frac{\alpha_2 \lambda_1}{\alpha_1\lambda_2}}\frac{b}{a} - 1.
\end{equation} 
This implies that 
\begin{equation}\label{eq:BoundThm4}
\left(\frac{\sinh\left(\left(\sqrt{\frac{\alpha_2}{\lambda_2}} b - \sqrt{\frac{\alpha_1}{\lambda_1}} a \right) \sqrt{r} \right)}{\sinh\left(\sqrt{\frac{\alpha_1}{\lambda_1}} a \sqrt{r}\right)\cosh\left(\sqrt{\frac{\alpha_2}{\lambda_2}} b \sqrt{r}\right)}\right)^k
\leq 2^k \left(\sqrt{\frac{\alpha_2 \lambda_1}{\alpha_1\lambda_2}}\frac{b}{a} - 1\right)^k \exp\left(-2k\sqrt{\frac{\alpha_1r}{\lambda_1}} a\right).
\end{equation} 
Note that the exponential part is slightly different compared with~\eqref{eq:BoundThm3} in the previous case. Using the bound~\eqref{eq:BoundThm4}, our estimate~\eqref{eq:min-cont} in this case becomes
\begin{equation}\label{eq:minBoundThm3}
\begin{aligned}
\norm{h^k}_{L^2(0, T)} \leq& \left(\frac{1}{1+\sqrt{\frac{\alpha_2\lambda_2}{\alpha_1\lambda_1}}}\right)^k\min_{r>0}e^{rT} \left(\hat H(r)\right)^k \norm{h^0}_{L^2(0, T)}
\\
\leq& \left(2\frac{\sqrt{\frac{\alpha_2 \lambda_1}{\alpha_1\lambda_2}}\frac{b}{a} - 1}{\sqrt{\frac{\alpha_2 \lambda_2}{\alpha_1\lambda_1}}+1}  \right)^k \min_{r>0}e^{rT - 2ka\sqrt{\frac{\alpha_1r}{\lambda_1}}}\norm{h^0}_{L^2(0, T)}.
\end{aligned}
\end{equation}
Once again, we take the derivative with respect to $r$ of the function $\exp(rT - 2ka\sqrt{\alpha_1r}/\sqrt{\lambda_1})$ to find the minimum 
\begin{equation}\label{eq:ropt-case2}
r^*_{2c} = \frac{\alpha_1 k^2a^2}{\lambda_1 T^2}.
\end{equation}
Substituting $r^*_{2c}$ back into the second inequality of~\eqref{eq:minBoundThm3} gives the estimate~\eqref{eq:estimate-case2-2}.
\end{proof}


For $\theta = \theta^*$, the error estimates given in Theorems~\ref{thm:cont-case1} and~\ref{thm:cont-case2} contain both linear and superlinear behavior. When the time $T$ is large, the exponential part in~\eqref{eq:ContEst1}, \eqref{eq:estimate-case2-1}, \eqref{eq:estimate-case2-2} is close to one, and the estimates predict linear convergence. In contrast, the exponential part plays an important role when the time $T$ is small, and we obtain superlinear convergence. 
Furthermore, the weight parameters~\eqref{eq:ropt-case1} and~\eqref{eq:ropt-case2} depend on the subdomain size and physical parameters: $r^*_{1c}$ depends on those in $Q_2$ and $r^*_{2c}$ depends on those in $Q_1$. 


On the other hand, the weight parameter $r$ is an auxiliary parameter that is only used in our analysis to derive error estimates. As shown in the proofs of Theorems~\ref{thm:cont-case1} and~\ref{thm:cont-case2}, our approach to treat the minimization problem~\eqref{eq:min-cont} is as follows: we first bound $\hat H(r)$ by a function of $r$, and then solve the resulting minimization problem to obtain $r^*_{1c}$ and $r^*_{2c}$, instead of solving~\eqref{eq:min-cont} directly. We can then substitute $r^*_{1c}$ and $r^*_{2c}$ into the function $\hat H(r)$ in~\eqref{eq:min-cont} to obtain sharper error estimates.

\begin{corollary}\label{thm:min-cont-bis}
For $\theta = \theta^*$, the error of the algorithm~\eqref{eq:DNWR} satisfies 
\begin{equation}\label{eq:min-cont-bis}
\norm{h^k}_{L^2(0, T)} 
\leq \left(\frac{ \hat H(r^*)}{1+\sqrt{\frac{\alpha_2\lambda_2}{\alpha_1\lambda_1}}}\right)^k e^{r^*T}\norm{h^0}_{L^2(0, T)},
\end{equation} 
where $r^*$ denotes $r_{1c}^*$ or $r_{2c}^*$ depending on the value of $\sqrt{\alpha_1/\lambda_1} a$ and $\sqrt{\alpha_2/\lambda_2} b$. 
\end{corollary}

Although the error estimate~\eqref{eq:min-cont-bis} is sharper than those in Theorems~\ref{thm:cont-case1} and~\ref{thm:cont-case2}, it is less straightforward to see the linear and superlinear convergence behavior in~\eqref{eq:min-cont-bis}, as the dependence on $T$ is hidden in the value of $r^*$. 

In general, one splits a long simulation time into smaller time windows, performs waveform relaxation methods within each window until a stopping criterion is satisfied, and then moves to the next one. Therefore, we are interested in the values of $T$ that lead to superlinear convergence. The error estimates provided in Theorems~\ref{thm:cont-case1} and~\ref{thm:cont-case2} indicate how one should choose the time window length, which will be further discussed in Section~\ref{sec:4}.
\section{Time discrete analysis}\label{sec:3}

The error estimates obtained in the continuous setting do not clarify the influence of the time step size. In this section, we provide a detailed time discrete analysis of~\eqref{eq:DNWR}. We discretize the time interval $[0, T]$ with $0=t_0 < t_1 < \ldots < t_N = T$, where $t_n = n\Delta t$ and $\Delta t = T/N$. We denote by $u^n$ the approximation of the solution $u$ at the time $t_n$ and $\boldsymbol{u} = (u^1, \ldots, u^N)^\top$. As before, we use $k$ to denote the index of the waveform iteration. The implicit Euler method applied to~\eqref{eq:DNWR} reads: For $n=0, \ldots, N-1$,
\begin{equation}\label{eq:DNWRIE}
	\begin{aligned}
		\frac{\alpha_1}{\Delta t} (u_1^{n+1, k+1} - u_1^{n, k+1}) - \lambda_1 \partial_{xx} u_1^{n+1, k+1} &= f_1^{n+1} \ \text{ in } \ (-a,0) , \\
		u_1^{n+1, k+1}(-a) = g_l^{n+1}, \quad 
		u_1^{n+1, k+1}(0) &= h^{n+1, k},\\
		u_1^{0, k+1}(x) &= u_{0, 1}(x) \\
		\frac{\alpha_2}{\Delta t} (u_2^{n+1, k+1}-u_2^{n, k+1})-\lambda_2 \partial_{xx} u_2^{n+1, k+1} &= f_2^{n+1} \ \text{ in } \ (0,b),\\
		u_2^{n+1,k+1}(b) = g_r^{n+1}, \quad
		-\lambda_1 \partial_x u_1^{n+1, k+1}(0) &= -\lambda_2 \partial_x u_2^{n+1, k+1}(0), \\
		u_2^{0, k+1}(x) &= u_{0, 2}(x),\\
		h^{n+1, k+1} := (1-\theta) h^{n+1, k} + \theta &u^{n+1, k+1}_2(0),  \quad \theta\in(0, 1).
	\end{aligned}
	\tag{DNWR-IE}
\end{equation}


To remain in a similar analysis framework as in Section~\ref{sec:2}, we keep using the exponentially weighted technique. In the continuous case, we use the Fourier transform; correspondingly, we consider here the discrete-time Fourier transform, which possesses properties, such as linearity, time shifting, differentiation, and the Parseval Theorem. We refer to the monograph~\cite[Chapter 2.9]{Oppenheim1999} for a detailed introduction. To take advantage of the Parseval Theorem, we define the exponentially weighted discrete-time Fourier transform as follows.


\begin{definition}\label{def:DiscFourier}
Let $r>0$, $\Delta t>0$, $\omega \in [-\pi, \pi]$,  and $\{\phi^n\}_{n\in\mathbb{Z}}$ be a given sequence. 
The exponentially weighted discrete-time Fourier transform of $\{\phi^n\}_{n\in\mathbb{Z}}$ is defined by
\[\widehat{\phi}(r\Delta t + i \omega) := \sum_{n\in\mathbb{Z}} \phi^n e^{-n(r\Delta t+i\omega)},\]
if this series is finite. 
\end{definition}


We use the notation wide hat here to denote the exponentially weighted discrete-time Fourier transform to distinguish it from the exponentially weighted Fourier transform in the continuous case. In addition, similar to Definition~\ref{def:FourierTransform}, we write the argument as $r\Delta t + i\omega$ to emphasize the dependence on the weight parameter $r$, the time step size $\Delta t$ and the frequency $\omega$. 
In general, the weight $e^{-rn\Delta t}$ blows up as $n\to -\infty$ for $r>0$. The next Lemma provides a similar result to Lemma~\ref{prepCont2norm}.

\begin{lemma}\label{prepDisc2norm}
Let $r > 0$ and $\boldsymbol{\phi} = \{\phi^n\}_{n\in\mathbb{Z}}$ be a sequence such that $ \sum_{n\in\mathbb{N}} (\phi^n)^2 < \infty$, and $\phi^n=0$ for $n<0$. Then the exponentially weighted discrete-time Fourier transform exists, and we have
\[\frac{1}{2 \pi} \int_{-\pi}^{\pi} |\widehat{\phi}(r\Delta t+i\omega)|^2 d\omega \leq \norm{\boldsymbol{\phi}}_{l^2}^2.\]
\end{lemma}


\begin{proof}
For $r>0$ and $\Delta t> 0$, $e^{-r \Delta t} < 1$. 
We then have $\sum_{n\in\mathbb{Z}}(e^{-r \Delta t} \phi^n)^2 \leq \sum_{n\in\mathbb{Z}} (\phi^n)^2 < \infty$. The time-discrete Fourier transform of the sequence $e^{-r \Delta t } \boldsymbol{\phi}$ therefore exists. Applying Parseval's theorem to $e^{-r \Delta t } \boldsymbol{\phi}$, we find
\[\frac{1}{2\pi} \int_{-\pi}^{\pi} \left|\widehat{\phi}(r \Delta t + i\omega)\right|^2 d \omega= \sum_{n\in\mathbb{Z}} (e^{-r \Delta t n} \phi^{n})^2 \leq \sum_{n\in\mathbb{Z}} (\phi^{n})^2 = \norm{\boldsymbol{\phi}}_{l^2}^2.\]
\end{proof}


To analyze convergence in the time discrete case, we focus on the error equation associated with~\eqref{eq:DNWRIE}, which consists of setting $f_1^n, f_2^n, g_l^n, g_r^n, u_{0, 1}, u_{0, 2}$ all to zero in~\eqref{eq:DNWRIE}. We also extend the initial guess $\{h^{n, 0}\}_{n=0}^{N}$ by 0 to $\boldsymbol{h}^0 :=\{h^{n,0}\}_{n\in\mathbb{Z}}$. The iterates $\{u_1^{n, k}, u_2^{n, k}, h^{n, k}\}_{n=0}^N$ for $k>0$ are then extended for $n\in\mathbb{Z}$ through the iteration itself. For brevity, we retain the notation $h^{n, 0}$ and $u_1^{n, k}, u_2^{n, k}, h^{n, k}$ for the extended version. 
The exponentially weighted discrete-time Fourier transform applied to the error equations gives
\begin{equation}\label{eq:transformedDNWRIE}
\begin{aligned}
	\frac{\alpha_1}{\Delta t} (e^{r \Delta t + i \omega} \widehat u^{k+1}_1 - \widehat u^{k+1}_1) - \lambda_1  \partial_{xx} e^{r \Delta t + i \omega} \widehat u^{k+1}_1 &= 0, \\
	\widehat u^{k+1}_1(-a) = 0, \quad
	\widehat u^{k+1}_1(0) &= \widehat h^k, \\
	\frac{\alpha_2}{\Delta t} (e^{r \Delta t + i \omega} \widehat u_2^{k+1} - \widehat u_2^{k+1}) - \lambda_2  \partial_{xx} e^{r \Delta t + i \omega} \widehat u_2^{k+1} &= 0, \\
	\widehat u^{k+1}_2(b) = 0, \quad
	-\lambda_1 \partial_x \widehat u^{k+1}_1(0) &= -\lambda_2 \partial_x \widehat u^{k+1}_2(0), \\
	\widehat h^{k+1} = (1 - \theta)\,\widehat h^{k} + &\theta\,\widehat u^{k+1}_2(0).
\end{aligned}
\end{equation}
Similar to~\eqref{eq:transformedDNWR}, these are two coupled second-order ordinary differential equations in space, and solving them separately yields the closed form solutions
\[\begin{aligned}
	\widehat u_1^k(x) &= \frac{\sinh\left((x+a)\sqrt{\frac{\alpha_1(1-e^{-r \Delta t - i\omega})}{\lambda_1\Delta t}}\right)}{\sinh\left(a\sqrt{\frac{\alpha_1(1-e^{-r \Delta t -i \omega})}{\lambda_1\Delta t}}\right)} \widehat h^k, \\
	\widehat u_2^k(x) &= \sqrt{\frac{\alpha_1\lambda_1}{\alpha_2\lambda_2}}\frac{\cosh\left(a\sqrt{\frac{\alpha_1(1-e^{-r \Delta t -i \omega})}{\lambda_1\Delta t}}\right)}{\sinh\left(a\sqrt{\frac{\alpha_1(1-e^{-r \Delta t -i \omega})}{\lambda_1\Delta t}}\right)}\frac{\sinh\left((x-b)\sqrt{\frac{\alpha_2(1-e^{-r \Delta t -i \omega})}{\lambda_2\Delta t}}\right)}{\cosh\left(b\sqrt{\frac{\alpha_2(1-e^{-r \Delta t -i \omega})}{\lambda_2\Delta t}}\right)} \widehat h^k.
\end{aligned}\]
Using the last equation in~\eqref{eq:transformedDNWRIE}, we obtain by induction that
\[\widehat h^k = \hat{\rho}\left(\frac{1-e^{-r \Delta t -i \omega}}{\Delta t}\right)^k \widehat h^0,\]
where the function $\hat{\rho}$ is defined in~\eqref{eq:cvfCont}. Similar to the continuous case, the convergence factor $\hat{\rho}((1-e^{-r \Delta t -i \omega})/\Delta t)$ depends on the subdomain sizes and the physical parameters that are all hidden in the function $\hat{\rho}$. Besides, the convergence factor also depends on the time step size $\Delta t$. Our goal is to understand how the choice of time step size influences the convergence behavior, and the relationship between the continuous and time discrete analysis. As in Theorem~\ref{2normSuperCont}, we first present an error estimate of the time discrete iteration~\eqref{eq:DNWRIE}
for arbitrary relaxation parameter $\theta\in(0, 1)$. 

\begin{theorem}\label{superLinearTimeDisc}
For any given $\theta\in(0, 1)$, the error of the time discrete iteration \eqref{eq:DNWRIE} satisfies the estimate 
\begin{equation}\label{eq:GenEstDis}
\|\boldsymbol{h}^k\|_{l^2} \leq \min_{r > 0} e^{rT}\max_{\omega \in [-\pi,\pi]} \left|\hat{\rho}\left(\frac{1-e^{-r \Delta t -i \omega}}{\Delta t}\right)^k\right| \|\boldsymbol{h}^0\|_{l^2},
\end{equation}
where $\hat{\rho}$ is defined in~\eqref{eq:cvfCont}.
\end{theorem}


\begin{proof}
As $r>0$, we have
\begin{align*}
	\|\boldsymbol{h}^k\|_{l^2}^2  = \sum_{n=0}^{N} e^{2 r \Delta t n}e^{-2 r \Delta t n} \left(h^{n,k}\right)^2 
	&\leq e^{2 rT} \sum_{n=0}^{N}  \left(e^{-r \Delta t n}h^{n,k}\right)^2 \\
	&\leq e^{2 rT} \sum_{n\in\mathbb{Z}} \left( e^{- r \Delta t n}h^{n,k}\right)^2,
\end{align*}
where in the last inequality we use the extension of the sequence $\boldsymbol{h}^k$. By Lemma~\ref{prepDisc2norm}, the exponentially weighted discrete-time Fourier transform of $\boldsymbol{h}^k$ exists. Applying the Parseval theorem yields 
\[\begin{aligned}
e^{2rT} \sum_{n\in\mathbb{Z}} \left( e^{- r \Delta t n}h^{n,k}\right)^2 
&= \frac{e^{2rT}}{2 \pi} \int_{-\pi}^{\pi} \left| \widehat{h}^{k}(r \Delta t + i \omega) \right|^2 d \omega \\
&= \frac{e^{2rT}}{2 \pi} \int_{-\pi}^{\pi} \left| \hat{\rho}\left(\frac{1-e^{-r \Delta t -i \omega}}{\Delta t}\right)^k \widehat{h}^{0}(r \Delta t + i \omega) \right|^2 d \omega \\
&\leq \frac{e^{2rT}}{2 \pi} \max_{\omega \in [-\pi,\pi]} \left|\hat{\rho}\left(\frac{1-e^{-r \Delta t -i \omega}}{\Delta t}\right)\right|^{2k}\int_{-\pi}^{\pi} \left| \widehat{h}^{0}(r \Delta t + i \omega) \right|^2  d \omega\\
&\leq e^{2rT} \max_{\omega \in [-\pi,\pi]} \left|\hat{\rho}\left(\frac{1-e^{-r \Delta t -i \omega}}{\Delta t}\right)\right|^{2k} \|\boldsymbol{h}^0\|_{l^2}^2,
\end{aligned}\]
where we use Lemma~\ref{prepDisc2norm} for the last inequality. Combining these computations gives
\[\|\boldsymbol{h}^k\|_{l^2}^2 \leq e^{2rT} \max_{\omega \in [-\pi,\pi]} \left|\hat{\rho}\left(\frac{1-e^{-r \Delta t -i \omega}}{\Delta t}\right)\right|^{2k} \|\boldsymbol{h}^0\|_{l^2}^2.\]
Taking the square root on both sides and minimizing over all $r > 0$ leads to~\eqref{eq:GenEstDis}.
\end{proof}


This estimate involves a min-max problem which is generally difficult to analyze. However, analogous to the continuous case, when the physical parameters in both subdomains satisfy $\sqrt{\alpha_1/\lambda_1}a = \sqrt{\alpha_2/\lambda_2}b$, the optimal relaxation parameter $\theta^*$ given in~\eqref{eq:thetaopt} again ensures convergence in two iterations for the time discrete iteration. Furthermore, we have obtained superlinear convergence in the continuous analysis using such $\theta^*$ for more general physical conditions, {\it i.e.}, $\sqrt{\alpha_1/\lambda_1}a \neq \sqrt{\alpha_2/\lambda_2}b$. We are interested in the convergence behavior of the time discrete iteration~\eqref{eq:DNWRIE} for this special choice of relaxation parameter $\theta^*$ and the dependence on the time step size $\Delta t$.
As in the continuous case, with the help of~\eqref{eq:rhotheta} and~\eqref{def:H}, we can rewrite the error estimate~\eqref{eq:GenEstDis} for $\theta=\theta^*$ as
\begin{equation}\label{eq:SpeEstDis}
\|\boldsymbol{h}^k\|_{l^2} \leq \left(\frac{1}{1+\sqrt{\frac{\alpha_2\lambda_2}{\alpha_1\lambda_1}}}\right)^k \min_{r > 0} e^{rT}\max_{\omega \in [-\pi,\pi]} \left|\hat H\left(\frac{1-e^{-r \Delta t -i \omega}}{\Delta t}\right)^k\right| \|\boldsymbol{h}^0\|_{l^2},
\end{equation} 
where $\alpha = \sqrt{\alpha_1/\lambda_1} a$ and $\beta = \sqrt{\alpha_2/\lambda_2} b$ in the function $\hat H$ defined in~\eqref{def:H}. Therefore, when $\theta=\theta^*$, it suffices to study the min-max problem associated with $\hat H$. We start by giving a similar result to Lemma~\ref{lem:max} for the global maximum of the maximization problem. 

\begin{lemma}\label{lem:maxDisc}
Let $\alpha> 0$, $\beta>0$, $r>0$, and $\Delta t>0$, then the function $\hat{H}$ defined in~\eqref{def:H} satisfies 
\[\max_{\omega \in [-\pi,\pi]} \left| \hat{H} \left(\frac{1-e^{-r \Delta t - i \omega}}{\Delta t}\right) \right| 
= \hat{H} \left(\frac{1-e^{-r \Delta t}}{\Delta t}\right).\]
\end{lemma}


\begin{proof}
Let $g(\omega) =(1-e^{-r \Delta t - i \omega})/ \Delta t$. For $r>0$ and $\Delta t>0$, the image of $g$ is a subset of a shifted half-plane, or more precisely
\[\text{Image}(g) = \left\{g(\omega): \, \omega \in [-\pi,\pi]\right\} \subset \left\{g(0) + x + iy \, : \,x\geq 0, y \in \mathbb{R} \right\} =: S.\]
Thus, we can transfer the maximization problem to the set $S$ as
\[\max_{\omega \in [-\pi,\pi]} \left| \hat{H}\left(g(\omega)\right) \right| \leq \max_{(x,y) \in S} \left| \hat{H} \left(g(0) + x + iy\right) \right| 
=  \max_{x \geq 0} \left| \hat H \left(g(0) + x\right) \right|,\]
where Lemma~\ref{lem:max} was used in the last equality. We now show that $\hat H$ is nonincreasing as a function of a real variable. Let $\gamma$ be a real number. Depending on the value of $\alpha$ and $\beta$, we have three cases.

\begin{enumerate}
\item[(i)] If $\alpha=\beta$, $\hat H = 0$. 

\item[(ii)] If $\alpha>\beta>0$, we have
\[\hat{H}'(\gamma) = \frac{\alpha \sinh(2 \beta \gamma) - \beta \sinh(2\alpha \gamma)}{2 \sinh^2(\alpha \gamma) \cosh^2(\beta \gamma)} < 0,\]
since $\alpha \sinh(2 \beta \gamma) < \beta \sinh(2\alpha \gamma)$. In fact, the function $f(\gamma) := \alpha \sinh(2 \beta \gamma) - \beta \sinh(2\alpha \gamma)$ satisfies
\[f(0) = 0, \quad f'(\gamma) = 2\alpha\beta(\cosh(2 \beta \gamma)-\cosh(2 \alpha \gamma)) < 0,\] 
for $\alpha > \beta >0$. 

\item[(iii)] If $\beta>\alpha>0$, we have
\[\hat{H}'(\gamma) = \frac{\beta \sinh(2 \alpha \gamma) -\alpha \sinh(2\beta \gamma)}{2 \sinh^2(a \gamma) \cosh^2(b \gamma)} < 0,\]
following a similar reasoning as in (ii). 
\end{enumerate}
Therefore, we have
\[\hat{H} \left(g(0) + x\right) \leq  \hat{H} \left(g(0)\right),
\quad \forall x \geq 0.\] 
This implies that
\[\max_{\omega \in [-\pi,\pi]} \left| \hat{H}\left(g(\omega)\right) \right|	
\leq \hat{H}\left(g(0)\right) = \hat{H}\left( \frac{1-e^{-r\Delta t}}{\Delta t} \right).\] 
\end{proof}


Applying Lemma~\ref{lem:maxDisc} to~\eqref{eq:SpeEstDis}, we have 
\begin{equation}\label{eq:min-disc}
\|\boldsymbol{h}^k\|_{l^2} \leq \left(\frac{1}{1+\sqrt{\frac{\alpha_2\lambda_2}{\alpha_1\lambda_1}}}\right)^k \min_{r > 0} e^{rT} \hat H\left(\frac{1-e^{-r \Delta t}}{\Delta t}\right)^k \|\boldsymbol{h}^0\|_{l^2}.
\end{equation}
As in the continuous case, we address the minimization problem for $r$ according to the underlying physics of the problem.


\begin{theorem}[Case $\sqrt{\alpha_1/\lambda_1} a > \sqrt{\alpha_2/\lambda_2} b$]\label{thm:timeDisc-case1}

If $\theta = \theta^*$ and $\sqrt{\alpha_1/\lambda_1} a > \sqrt{\alpha_2/\lambda_2} b$, then the error of the time discrete iterates~\eqref{eq:DNWRIE} satisfies
\begin{equation}\label{eq:DiscEst1}
\|\boldsymbol{h}^k\|_{l^2} \leq  \left(2\frac{1 - \sqrt{\frac{\alpha_2 \lambda_1}{\alpha_1\lambda_2}}\frac{b}{a}}{1 + \sqrt{\frac{\alpha_2 \lambda_2}{\alpha_1\lambda_1}}} \right)^k \phi_1(r_{1d}^*) \norm{\boldsymbol{h}^0}_{l^2},
\end{equation}
with
\begin{equation}\label{eq:roptDisc-case1}
\phi_1(r) = \exp\left( rT  - kb\sqrt{\frac{\alpha_2}{\lambda_2}}\sqrt{\frac{1-e^{-r \Delta t}}{\Delta t}}\right), \ r^*_{1d} = \frac1{\Delta t}\log\left(\frac{1 + \sqrt{\frac{\alpha_2\Delta t}{\lambda_2T^2} b^2 k^2   + 1}}{2}\right).
\end{equation}
\end{theorem}

\begin{proof}
Reusing the bound~\eqref{eq:BoundThm2} given in the proof of Theorem~\ref{thm:cont-case1} with $r = (1-\exp(-r \Delta t))/\Delta t$ yields
\begin{equation}\label{eq:minBoundThm6}
\left(\frac1{1+\sqrt{\frac{\alpha_2\lambda_2}{\alpha_1\lambda_1}}}\right)^k\min_{r>0}e^{rT} \hat H \left(\frac{1-e^{-r \Delta t}}{\Delta t}\right)^k 
\leq \left(2\frac{\bigl|1 - \sqrt{\frac{\alpha_2 \lambda_1}{\alpha_1\lambda_2}}\frac{b}{a}\bigr|}{1 + \sqrt{\frac{\alpha_2 \lambda_2}{\alpha_1 \lambda_1}}} \right)^k\min_{r > 0} \phi_1(r).
\end{equation}
 Differentiating $\phi_1$ yields
\[\phi_1'(r) = \left(T - \frac{kb}{2}\sqrt{\frac{\alpha_2}{\lambda_2}}\sqrt{\frac{\Delta t}{1 - e^{-r \Delta t}}} e^{-r \Delta t} \right)\exp\left(r T -  kb\sqrt{\frac{\alpha_2}{\lambda_2}} \sqrt{\frac{1 - e^{-r \Delta t}}{\Delta t}} \right).\]
Setting the latter to 0 yields the extremal point $r^*_{1d}$ given in~\eqref{eq:roptDisc-case1}. The derivative $\phi_1'(r)$ is negative for $r<r^*_{1d}$ and positive for $r>r^*_{1d}$, since $e^{-r \Delta t}\sqrt{\Delta t/(1 - e^{-r \Delta t})}$ is a decreasing function of $r$. This implies that the extremum is indeed a minimum. Inserting $r^*_{1d}$ into $\phi_1$ in~\eqref{eq:minBoundThm6} gives the estimate~\eqref{eq:DiscEst1}.
\end{proof}



\begin{theorem}[Case $\sqrt{\alpha_1/\lambda_1} a < \sqrt{\alpha_2/\lambda_2} b$]\label{thm:timeDisc-case2}

Let $\theta = \theta^*$. If $\sqrt{\alpha_1/\lambda_1} a < \sqrt{\alpha_2/\lambda_2} b < 2\sqrt{\alpha_1/\lambda_1} a$, then the error of the iterates~\eqref{eq:DNWRIE} satisfies
\begin{equation}\label{eq:DiscEst2-1}
\|\boldsymbol{h}^k\|_{l^2} \leq \left(2\frac{\sqrt{\frac{\alpha_2 \lambda_1}{\alpha_1\lambda_2}}\frac{b}{a} -1}{\sqrt{\frac{\alpha_2 \lambda_2}{\alpha_1\lambda_1}} + 1} \right)^k \phi_1(r^*_{1d}) \|\boldsymbol{h}^0\|_{l^2},
\end{equation}
with $r^*_{1d}$ and $\phi_1$ defined in~\eqref{eq:roptDisc-case1}. 

If $\sqrt{\alpha_2/\lambda_2} b \geq 2\sqrt{\alpha_1/\lambda_1} a$, then the error of the iterates~\eqref{eq:DNWRIE} satisfies
\begin{equation}\label{eq:DiscEst2-2}
\|\boldsymbol{h}^k\|_{l^2} \leq \left(2\frac{\sqrt{\frac{\alpha_2 \lambda_1}{\alpha_1\lambda_2}}\frac{b}{a} - 1}{\sqrt{\frac{\alpha_2 \lambda_2}{\alpha_1\lambda_1}}+1} \right)^k 
\phi_2(r_{2d}^*) \|\boldsymbol{h}^0\|_{l^2},
\end{equation}
with 
\begin{equation}\label{eq:roptDisc-case2}
\phi_2(r) = \exp\left( r T  - ka\sqrt{\frac{\alpha_1}{\lambda_1}}\sqrt{\frac{1-e^{-r \Delta t}}{\Delta t}}\right), \  r^*_{2d} = \frac1{\Delta t}\log\left(\frac{1 + \sqrt{4\frac{\alpha_1  \Delta t}{ \lambda_1 T^2}a^2 k^2 + 1}}{2}\right).
\end{equation}
\end{theorem}

\begin{proof}
The proof for the case $\sqrt{\alpha_1/\lambda_1} a < \sqrt{\alpha_2/\lambda_2} b < 2\sqrt{\alpha_1/\lambda_1} a$ follows the same steps as that of Theorem~\ref{thm:timeDisc-case1}, except we make use of the bound~\eqref{eq:BoundThm3}. This leads to the same optimized $r^*_{1d}$ given in~\eqref{eq:roptDisc-case1}.
	

For $\sqrt{\alpha_2/\lambda_2} b \geq 2\sqrt{\alpha_1/\lambda_1} a$, we follow once again the proof of Theorem~\ref{thm:timeDisc-case1}, but use instead  the bound~\eqref{eq:bound-exp}, which gives
\[\left(\frac1{1+\sqrt{\frac{\alpha_2\lambda_2}{\alpha_1\lambda_1}}}\right)^k\min_{r>0}e^{rT} \hat H \left(\frac{1-e^{-r \Delta t}}{\Delta t}\right)^k 
\leq \left(2\frac{\frac{\sqrt{\alpha_2 \lambda_1}b}{\sqrt{\alpha_1\lambda_2}a} - 1}{\sqrt{\frac{\alpha_2 \lambda_2}{\alpha_1\lambda_1}}+1} \right)^k \min_{r > 0} \phi_2(r).\]
Setting then $\phi'_2(r)=0$, we obtain the optimum value of $r^*_{2d}$~\eqref{eq:roptDisc-case2}. Inserting it into $\phi_2$ gives the result~\eqref{eq:DiscEst2-2}.
\end{proof}



Theorems~\ref{thm:timeDisc-case1} and~\ref{thm:timeDisc-case2} provide time discrete error estimates that depend on both time $T$ and step size $\Delta t$. Note that $\phi_1(r^*_{1d})$ and $\phi_2(r^*_{2d})$ are guaranteed to be less than or equal to 1, since $\phi_1$ and $\phi_2$ are continuous for $r>0$ and $1 =\phi_1(0) \geq \min_{r>0} \phi_1(r) = \phi_1(r^*_{1d})$ and $1 =\phi_2(0) \geq \min_{r>0} \phi_2(r) = \phi_2(r^*_{2d})$. Thus, the exponential factor is guaranteed to be nonincreasing. 

Furthermore, the time discrete estimates share the same linear factor as in the continuous estimates given in Theorems~\ref{thm:cont-case1} and~\ref{thm:cont-case2}. Similarly to the continuous case, the weights~\eqref{eq:roptDisc-case1} and~\eqref{eq:roptDisc-case2} depend on the subdomain size and physical parameters. As they are auxiliary parameters that are only used in our analysis to derive error estimates, we can also substitute them directly into~\eqref{eq:min-cont} to obtain sharper error estimates, as in Corollary~\ref{thm:min-cont-bis}.


\begin{corollary}\label{thm:min-disc-bis}
	For $\theta = \theta^*$, the error of the iterates~\eqref{eq:DNWRIE} satisfies
	\begin{equation}\label{eq:min-disc-bis}
			\|\boldsymbol h^k\|_{l^2} 
			\leq \left(\frac{\hat H \left(\frac{1-e^{-r^* \Delta t}}{\Delta t}\right)}{1+\sqrt{\frac{\alpha_2\lambda_2}{\alpha_1\lambda_1}}}\right)^ke^{r^*T} \|\boldsymbol{h}^0\|_{l^2},
	\end{equation} 
	where $r^*$ denotes $r_{1d^*}$ or $r_{2d^*}$ depending on the value of  $\sqrt{\alpha_1/\lambda_1} a$ and $\sqrt{\alpha_2/\lambda_2} b$.
\end{corollary}


Once again, the error estimate~\eqref{eq:min-disc-bis} is sharper than those in Theorems~\ref{thm:timeDisc-case1} and~\ref{thm:timeDisc-case2}, since we directly inject the value $r^*$ (i.e., $r_{1d}^*$ or $r_{2d}^*$) into the function $\hat H$ instead of bounding it first. However, its dependence on the time step size $\Delta t$ is hidden in the value of $r^*$. To compare with the estimates obtained in the continuous case and to understand the impact of $\Delta t$, we still need to use the estimates derived in Theorems~\ref{thm:timeDisc-case1} and~\ref{thm:timeDisc-case2}.
\section{Comparison of continuous and time discrete estimates}\label{sec:4}


\begin{thebibliography}{99}
\bibitem[]{Ciaramella2022} Gabriele Ciaramella; Martin J. Gander. Iterative Methods and Preconditioners for Systems of Linear Equations. (2022)
\bibitem[]{Kwok2025} Alejandro Alfonso Rodriguez; Felix Kwok; Martin J. Gander. Convergence analysis of optimized {S}chwarz waveform relaxation by exponential weighting. Submitted (2026)
\bibitem[]{Lions72} Lions, J. L.; Magenes, E.. Non-Homogeneous Boundary Value Problems and Applications II. 182 (1972)
\bibitem[]{Oppenheim1999} Alan V. Oppenheim; Ronald W. Schafer. Discrete-Time Signal Processing. (1999)
\bibitem[]{grafakos2014a} Grafakos, Loukas. Classical Fourier Analysis. 249 (2014)
\bibitem[]{ma:13} Mandal, Bankim C.. A Time-Dependent {D}irichlet-{N}eumann Method for the Heat Equation. Domain Decomposition Methods in Science and Engineering XXI 98 (2014) 467--475
\end{thebibliography}
\end{document}
